% Propietario conservado: /Users/ruben/Documents/New project/output/REVISION_ARGUMENTAL_APERTURAS_CIERRES_20260915/VII/delivery/MANUSCRIPT_VII_ES_EN_COMPLETE/source_es/manuscrito/30d_realizacion_geometrica.tex
% RESULTADO_RECUPERADO: S15, 11c_gravedad_holonomia_rev6.tex:128--212;
% S01, 02_dinamica_tres_hojas.tex:155--219; pantalla y respuesta de30.
% Pruebas de curvatura, covariancia y Bianchi reunidas para autonomía.
% La realización suave conserva sus condiciones de holonomía explícitas.

\section{Réalisation géométrique du transport et de la soudure}
\label{vii:vii:sec:realizacion-cartan}

La représentation discrète du retour et la soudure d'écran déterminent deux aspects du transport : l'accumulation de mémoire dans les routes et la comparaison de leurs fibres. Leur réalisation géométrique conserve les deux opérations. Dans cette section sont fixés l'espace d'arrivée, la normalisation dimensionnelle de la connexion et les identités qu'utilisera sa variation.

\subsection{Écran, cotétrade et connexion}

Le quotient d'écran \(Q_\square\), sa forme lorentzienne et la
soudure \(\beta_m\) procèdent de
\ref{vii:vii:sec:pantalla-respuesta}. Dans chaque cellule, une réalisation d'
écran \(J_{{\rm scr},m}\) transporte la forme et la soudure
vers la fibre géométrique correspondante. Son image cellulaire est
\[
 e_m=J_{{\rm scr},m}\beta_m.
\]
L'échelle \(\ell_m\) demeure à l'intérieur de la soudure. Les
changements de repère agissent sur \(J_{{\rm scr},m}\), la connexion
et le transport des cellules conjointement.

Une réalisation lisse de ces données sur une variété orientée \(\mathcal U\) de dimension quatre consiste en une cotétrade non dégénérée \(e=(e^I)\), une connexion de Lorentz \(\omega=(\omega^I{}_J)\) et une réalisation de routes \(\mathfrak R_{\rm C}\) qui satisfont
\begin{equation}
 \operatorname{Hol}_{\mathcal A_{\ell_0}}
       (\mathfrak R_{\rm C}(\gamma))
 =\rho_{\rm C}(\operatorname{Hol}_{\rm TPK}(\gamma)).
 \label{vii:vii:eq:compatibilidad-holonomia-geometrica}
\end{equation}
Dans cette égalité matricielle, la représentation est normalisée par
\[
 N_{\ell_0}(R,a)=(\Lambda(R),a/\ell_0),\qquad
 \rho_{\rm C}=N_{\ell_0}\circ\rho_{\rm C}^{\rm disc},
\]
où \(\Lambda:\operatorname{Spin}^+(3,1)\to SO^+(3,1)\)
est la représentation vectorielle. L'application \(N_{\ell_0}\)
préserve le produit semi-direct : diviser la translation
\(a+\Lambda(R)b\) par \(\ell_0\) équivaut à composer les
deux translations normalisées. Le tour complet publie donc
\(72u\) dans la connexion normalisée et \(72\ell_0u\) dans la
carte dimensionnelle. Le relèvement de spin conserve, séparément,
\(R_4^{c_g(\gamma)}\) ; la représentation vectorielle a pour
noyau central \(\{1,-1\}\), de sorte que ce
relèvement est retenu lors de l'utilisation de champs spinoriels.
L'identité, la concaténation, l'inversion et la subdivision sont maintenues dans la
réalisation correspondante. L'appariement d'
écran et la condition d'holonomie spécifient les données de la
réalisation utilisée dans les équations de champ.

Avec \(\eta=\operatorname{diag}(-1,1,1,1)\), on définit
\(g=\eta_{IJ}e^I\otimes e^J\). Pour une échelle élémentaire
positive \(\ell_0\), constante dans la carte, la connexion affine est
\begin{equation}
 \mathcal A_{\ell_0}
 =\begin{pmatrix}\omega&\ell_0^{-1}e\\0&0\end{pmatrix}.
 \label{vii:vii:eq:conexion-afin-realizada}
\end{equation}
La cotétrade a la dimension d'une longueur ; le facteur \(\ell_0^{-1}\) homogénéise sa composante translationnelle.

\begin{proposition}[Courbure et torsion de la connexion réalisée]
\label{vii:vii:prop:curvatura-afin}
La courbure de \eqref{vii:vii:eq:conexion-afin-realizada} est
\begin{equation}
 \mathcal F_{\ell_0}
 =\mathrm d\mathcal A_{\ell_0}
   +\mathcal A_{\ell_0}\wedge\mathcal A_{\ell_0}
 =\begin{pmatrix}F_\omega&\ell_0^{-1}T\\0&0\end{pmatrix},
 \quad F_\omega=\mathrm d\omega+\omega\wedge\omega,
 \quad T=\mathrm de+\omega\wedge e.
 \label{vii:vii:eq:curvatura-afin-realizada}
\end{equation}
\end{proposition}
\begin{proof}
Le produit de matrices de formes produit \(\omega\wedge\omega\) dans le bloc diagonal et \(\ell_0^{-1}\omega\wedge e\) dans le bloc translationnel. La dérivée extérieure apporte \(\mathrm d\omega\) et \(\ell_0^{-1}\mathrm de\), puisque \(\ell_0\) est constant. La somme démontre l'égalité.
\end{proof}

La courbure décrit le défaut de transport des repères et la torsion
celui de la cotétrade. Toutes deux appartiennent à la même connexion réalisée ;
leurs composantes conservent des domaines tensoriels distincts.

\begin{proposition}[Covariance et identités de Bianchi]
\label{vii:vii:prop:bianchi-cartan}
Sous un changement de repère \(g_L:\mathcal U\to SO^+(3,1)\),
\[
 e'=g_Le,\qquad
 \omega'=g_L\omega g_L^{-1}-\mathrm dg_L\,g_L^{-1},
\]
on a \(T'=g_LT\) et \(F_{\omega'}=g_LF_\omega g_L^{-1}\). En outre,
\begin{equation}
 D_\omega T=F_\omega\wedge e,\qquad D_\omega F_\omega=0.
 \label{vii:vii:eq:bianchi-realizacion}
\end{equation}
\end{proposition}
\begin{proof}
Dans \(\mathrm d(g_Le)+\omega'\wedge g_Le\), les termes
contenant \(\mathrm dg_L\wedge e\) se compensent.
La même substitution dans la courbure donne sa transformation par
conjugaison. Pour la première identité de Bianchi, on développe
\[
 D_\omega(\mathrm de+\omega\wedge e)
 =(\mathrm d\omega+\omega\wedge\omega)\wedge e.
\]
Dans \(\mathrm dF_\omega+\omega\wedge F_\omega-F_\omega\wedge\omega\), les termes contenant \(\mathrm d\omega\) et les produits cubiques s'annulent par paires. Cela prouve la seconde identité.
\end{proof}

\subsection{Les trois réalisations quadratiques de courbure}

Le feuillet \(\tau\in\{+1,0,-1\}\) du TRIT est représenté dans
l'algèbre de Cartan au moyen des générateurs \(J_{ab}=-J_{ba}\)
et \(P_a^{(\tau)}\), avec les relations
\begin{align*}
 [J_{ab},J_{cd}]&=
 \eta_{bc}J_{ad}-\eta_{ac}J_{bd}
 -\eta_{bd}J_{ac}+\eta_{ad}J_{bc},\\
 [J_{ab},P_c^{(\tau)}]&=
 \eta_{bc}P_a^{(\tau)}-\eta_{ac}P_b^{(\tau)},\\
 [P_a^{(\tau)},P_b^{(\tau)}]&=-\tau J_{ab}.
\end{align*}
Pour \(\ell>0\) constant, on écrit
\[
 \mathcal A_\tau^{\rm Cart}
 =\tfrac12\omega^{ab}J_{ab}+\ell^{-1}e^aP_a^{(\tau)}.
\]

\begin{proposition}[Courbure de la réalisation tritifiée]
\label{vii:vii:prop:curvatura-tritificada}
Avec les relations précédentes,
\begin{equation}
 \mathcal F_\tau^{\rm Cart}
 =\tfrac12(F_\omega^{ab}-\tau\ell^{-2}e^a\wedge e^b)J_{ab}
     +\ell^{-1}T^aP_a^{(\tau)}.
 \label{vii:vii:eq:curvatura-tritificada}
\end{equation}
\end{proposition}
\begin{proof}
On développe \(\mathrm d\mathcal A+\tfrac12[\mathcal A,\mathcal A]\). Le crochet du secteur de Lorentz produit \(F_\omega\), le crochet mixte produit \(D_\omega e\), et le crochet translationnel apporte \(-\tfrac12\tau\ell^{-2}e^a\wedge e^bJ_{ab}\).
\end{proof}

Dans le feuillet central, on retrouve la connexion affine ; les feuillets extrêmes
conservent les deux signes du terme de courbure constante. La
représentation maintient ainsi le régime tritique de la route avant
d'appliquer les équations variationnelles. L'égalité d'holonomie
\eqref{vii:vii:eq:compatibilidad-holonomia-geometrica} correspond au
feuillet central. Dans chaque feuillet extrême, on utilise la représentation
\(\rho_{{\rm C},\tau}\) dans son groupe de Cartan et sa condition
d'holonomie respective. Le courant et le tenseur métrique utilisent
les champs et la représentation de la même réalisation.
